\documentclass[output=paper]{langscibook}
\ChapterDOI{10.5281/zenodo.22880248}
\title{\textit{Kara naono} – Vines, aerial roots, fibres and strings}
\author{Ruth Sima Rigamu and Ondria Tavagaga}
\abstract{}
\begin{document}
\maketitle\label{chap:11}
\vspace{-1cm}
\begin{figure}[H]
    \captionsetup[subfigure]{oneside,margin={0cm,0cm}}
    \captionsetup{justification=centering}
    \subfloat[]{\includegraphics[trim=0 0 0 0,clip,width=.33\linewidth]{figures/11_1.jpg}\label{fig:11_1}}\hfil
    \subfloat[]{\includegraphics[trim=0 0 0 40,clip,width=.34\linewidth]{figures/11_2.jpg}\label{fig:11_2}}\hfil
    \subfloat[]{\includegraphics[trim=0 0 0 0,clip,width=.33\linewidth]{figures/11_3.jpg}\label{fig:11_3}}
    \caption[(a) \textit{Kara aisi} -- A weed, (b) \textit{Kara ariono} -- A vine, (c) \textit{Kara kanono} -- Some strings]{(a) \textit{Kara aisi} -- A weed, (b) \textit{Kara ariono} -- A vine,\\ (c) \textit{Kara kanono} -- Some strings\footnotemark}
\end{figure}\footnotetext{The word \textit{kara} is a kind of classifier for plants and other things that look like a string. In Adjectival Phrases that relate to a \textit{kara}-N compound as a predicate or modifier this classifier is repeated, see Chapter \ref{chap:11}:\ref{ex:text7-48} and Chapter \ref{chap:11}:\ref{ex:text7-78}.}
\begin{xltabular}{\textwidth}{lXX}
    (\ref{ex:text7-1})&\textbf{\textit{Kara aisi}}&\textbf{The weed called \textit{aisi}}\\
    (\ref{ex:text7-2})&\textit{O kara aisi, eve o meho kara naono to abeabe batana tea maa kasuana vaa Hanatana.}&The \textit{aisi}, it is a string-like plant that is creeping on the beaches of Hanatana.\\
    (\ref{ex:text7-3})&\textit{O kara aisi bona na tei nana kasuana.}&This \textit{aisi} is found on the ground.\\
    (\ref{ex:text7-4})&\textit{Ahiki vira ta kiu teve vo teara.}&Actually, it does not have any use for us.\\
\end{xltabular}
\begin{xltabular}{\textwidth}{lXX}
    (\ref{ex:text7-5})&\textit{\textbf{Kara ariono}}&\textbf{\textit{Ariono} vine}\\
    (\ref{ex:text7-6})&\textit{O kara ariono mee, eve o meho kara koris o kara mataa.}&The \textit{ariono} vine, it is a good vine.\\
    (\ref{ex:text7-7})&\textit{Eve to vagana ni ria abana tahii.}&It is used by the men for fishing in the sea.\\
    (\ref{ex:text7-8})&\textit{Beori nao tahii, eori toro gono bona kara ariono mene a vavaatoa.}&When they go (fishing) in the sea, they must get the \textit{ariono} string for kite-fishing.\footnote{See Chapter \ref{chap:34}:\ref{ex:text27-8}. When used as a name for the vine, \textit{kara ariono} is a N3-construction; but when used in the sense of `string made from \textit{ariono} vines' it is a N2-construction.}\\
    (\ref{ex:text7-9})&\textit{Na paku vavaatoa ni rori bona kara ariono.}&They do kite fishing with the \textit{ariono} string.\\
\end{xltabular}
\begin{xltabular}{\textwidth}{lXX}
	(\ref{ex:text7-10})&\textbf{\textit{Kara bao}}&\textbf{Grass}\\
	(\ref{ex:text7-11})&\textit{A kara bao eve a tabaan teo burumakao.}&Grass is the food of cows.\\
	(\ref{ex:text7-12})&\textit{Nabunuu roho enam ahiki to kara bao.}&In the past we did not have any grass.\\
	(\ref{ex:text7-13})&\textit{O burumakao he nomaa koa, meori paa ma neba bata kiri bari bono kara bao.}&But the cows came, and they planted the grass for them.\\
	(\ref{ex:text7-14})&\textit{Mea naba kasi hanana paa dee taataana bata maa bari.}&And the road builders brought it here along the coast.\\
	(\ref{ex:text7-15})&\textit{Rerau, mepaa raha bata maa.}&And so, it spread out here.\\
	(\ref{ex:text7-16})&\textit{Me nomana, amaa mohina teara na hapus kurus nana tea kara bao.}&And today our gardens are full of grass.\\
\end{xltabular}
\begin{xltabular}{\textwidth}{lXX}
	(\ref{ex:text7-17})&\textbf{\textit{Kara kave}}&\textbf{Net vine}\\
	(\ref{ex:text7-18})&\textit{O kara kave, eve o meho kara koris vai to tei me nana bona kiu a beera sana.}&The net vine is a vine that has a very important function.\\
	(\ref{ex:text7-19})&\textit{Nabunuu roho beo kara kave vai hiki nana,}&In the past, when there was not any net vine,\\
	(\ref{ex:text7-20})&\textit{ahiki koa ta an iana teara.}&we could not eat any fish.\\
	(\ref{ex:text7-21})&\textit{Eikou o kara kave to paku riori bona kanono,}&Because (it was) the net vine from which they made the string,\\
	(\ref{ex:text7-22})&\textit{ore paa paku pete bari bona kave mene vagana.}&and then they made the net for fishing from it.\footnote{See Chapter \ref{chap:29}:\ref{ex:text22-37}ff.}\\
\end{xltabular}
\begin{figure}[H]
    \captionsetup[subfigure]{oneside,margin={0cm,0cm}}
    \captionsetup{justification=centering}
    \subfloat[]{\includegraphics[trim=0 0 0 0,clip,width=.4\linewidth]{figures/11_4.png}\label{fig:11_4}}\hfil
    \subfloat[]{\includegraphics[trim=0 0 0 0,clip,width=.4\linewidth]{figures/11_5.jpg}\label{fig:11_5}}\hfil
    \subfloat[]{\includegraphics[trim=0 0 0 0,clip,width=.2\linewidth]{figures/11_6.png}\label{fig:11_6}}
    \caption{(a) \textit{Kara koris} -- Vine, (b) \textit{Kara kuskus} -- Passionflower,\\ (c) \textit{Obai} -- Bengal calamus}
\end{figure}
\begin{xltabular}{\textwidth}{lXX}
	(\ref{ex:text7-23})&\textbf{\textit{Kara koris}}&\textbf{Vine (generic term)}\\
	(\ref{ex:text7-24})&\textit{O kara koris, eve o meho kara naono.}&The vine is a string-like plant.\\
	(\ref{ex:text7-25})&\textit{O kara koris bona na tei pete nana rasuu ei mohina.}&The vine also grows in the jungle, in the garden.\\
	(\ref{ex:text7-26})&\textit{O kara koris bona na vaahiroo ni raara, na dao raara bono kara koremana.}&This vine we use for soup, we call \textit{koremana} vine.\\
	(\ref{ex:text7-27})&\textit{Bea moon vahuhu, ore paa vaahiroo bari bono kara koremana, ei bea huhu na moon vagitagita nana,}&When a woman gives birth, they give her the \textit{koremana} vine to drink as soup, if the woman's breasts are empty,\\
	(\ref{ex:text7-28})&\textit{a huhu nae repaa hapus.}&so that her breasts become full.  \\
	(\ref{ex:text7-29})&\textit{A pate koris bona o paapao.}&This bunch of vines is a vegetable.\\
	(\ref{ex:text7-30})&\textit{Eve o paapao o mataa.}&It is a good vegetable.\\
	(\ref{ex:text7-31})&\textit{``Hiroo o kara korisi bona.''}&``Drink this vine as a soup.'' \\
\end{xltabular}
\begin{xltabular}{\textwidth}{lXX}
	(\ref{ex:text7-32})&\textbf{\textit{Kara kuskus}}&\textbf{Passionflower}\\
	(\ref{ex:text7-33})&\textit{O kara kuskus eve mee o meho kara koris to tei nana mohina.}&The passionflower is a vine that is found in the garden.\\
	(\ref{ex:text7-34})&\textit{O kara kuskus bona na vuavua nana.}&The passionflower bears fruit.\\
	(\ref{ex:text7-35})&\textit{Bea vuanae manoto, ore paa tomtomon,}&When its fruit matures, it becomes orange,\\
	(\ref{ex:text7-36})&\textit{ei be vuri.}&this means that it is ripe.\\
	(\ref{ex:text7-37})&\textit{``Enaa na rake nom tea ani to vua kuskus.}&``I want to eat a passionflower fruit.\\
	(\ref{ex:text7-38})&\textit{O vua kuskus bona o vua mataa kurus.''}&This passionflower fruit is a very good fruit.''\\
	(\ref{ex:text7-39})&\textit{O kara kuskus na tei nana rasuu.}&The passionflower is found in the jungle.\\
	(\ref{ex:text7-40})&\textit{Na hiri batana tea maa pasunana tenam.}&It is creeping at the head of our gardens.\\
	(\ref{ex:text7-41})&\textit{Be vua, ore paa hana batana teo kara nae.}&When it bears fruit, it hangs on its vines.\\
	(\ref{ex:text7-42})&\textit{Beove vuri ore paa tomtomon.}&When it ripens, it becomes orange.\\
	(\ref{ex:text7-43})&\textit{A vua no kara kuskus a vua mataa kurus.}&The fruit of the passionflower is very good.\\
\end{xltabular}
\begin{xltabular}{\textwidth}{lXX}
	(\ref{ex:text7-44})&\textbf{\textit{Kara obai}}&\textbf{Calamus cane, Sweet Flag}\\
	(\ref{ex:text7-45})&\textit{O obai, a kiu nae a beera kurus.}&The calamus cane, its use is very important.\footnote{See Chapter \ref{chap:6}:\ref{ex:text2-32}ff and Chapter \ref{chap:17}:\ref{ex:text12-38}.}\\
	(\ref{ex:text7-46})&\textit{O kara obai, eve to paku riori bona sia bara na paku rori bona teevoro ae a teebana.}&The calamus cane is used for making chairs and also used to make tables and beds.\\
	(\ref{ex:text7-47})&\textit{O kara obai, eve mee to kisi ni riori bona inu bara tea vaave mee.}&The calamus cane is also used for lashing (the posts and beams of) the house and also for thatching (the roof).\\
	(\ref{ex:text7-48})&\textit{O kara obai kou o kara kikis kurus. Na hiri nana teo kasuana, ore paa abe hiava teo naono.}&The calamus cane is very strong. It creeps on the ground, and then it creeps up a tree.\\
	(\ref{ex:text7-49})&\textit{O paka obai na toutou nana.   }&The calamus leaf is prickly.  \\
\end{xltabular}
\begin{figure}[H]
    \captionsetup[subfigure]{oneside,margin={0cm,0cm}}
    \captionsetup{justification=centering}
    \subfloat[]{\includegraphics[trim=0 17 0 0,clip,width=.26\linewidth]{figures/11_7.jpg}\label{fig:11_7}}
    \subfloat[]{\includegraphics[trim=0 30 0 30,clip,width=.45\linewidth]{figures/11_8.jpg}\label{fig:11_8}}
    \subfloat[]{\includegraphics[trim=0 0 0 0,clip,width=.29\linewidth]{figures/11_9.jpg}\label{fig:11_9}}
    \caption{(a) \textit{Veresun} -- Wild banana, (b) \textit{Kara orovi} -- Black sheaths,\\ (c) \textit{Vin maata} -- Plaiting mats}
\end{figure}
\begin{xltabular}{\textwidth}{lXX}
	(\ref{ex:text7-50})&\textbf{\textit{Kara orovi}}&\textbf{Black sheath of the wild banana}\\
	(\ref{ex:text7-51})&\textit{O kara orovi na kurus maa raara teo veresun.}&The black wild banana sheath we cut off the wild banana plant.\\
	(\ref{ex:text7-52})&\textit{O veresun bona o paru na dao raara bono orovi.}&This wild banana plant (with) black (leaf sheaths), we call \textit{orovi}.\\
	(\ref{ex:text7-53})&\textit{Are paa dee voosu mae,}&(When we have cut off a black sheath,) we carry it home,  \\
	(\ref{ex:text7-54})&\textit{are paa ma kahu e bona kehaa.}&and scrape it with a seashell.\\
	(\ref{ex:text7-55})&\textit{Kahu vakavara, are paa varoava,}&After having scraped (it),\\
	(\ref{ex:text7-56})&\textit{are paa varoava e, ore raka.}&we put it in the sun so that it becomes dry.\\
	(\ref{ex:text7-57})&\textit{Be raka, are paa rigunue,}&When (it) is dry, we roll it up,\\
	(\ref{ex:text7-58})&\textit{matatopo vaa tea vin maata.}&in preparation for weaving mats.\\
	(\ref{ex:text7-59})&\textit{Are paa oroorovi nio orovi bona teo vonaha.}&Then we weave this black sheath into a pandanus mat.\footnote{See Chapter \ref{chap:13}:\ref{ex:text9-2}.}\\
\end{xltabular}
\begin{xltabular}{\textwidth}{lXX}
	(\ref{ex:text7-60})&\textbf{\textit{Kara Pihi}}&\textbf{Farting Weed}\\
	(\ref{ex:text7-61})&\textit{O Kara Pihi o meho subui vaa komana vaan.}&The Farting Weed is a weed in the village.\\
	(\ref{ex:text7-62})&\textit{O hena Kara Pihi na tei me nana bona ei to kao vo nae.}&The name Farting Weed has been given to it for the following reason.\\
	(\ref{ex:text7-63})&\textit{Bean kara e, saka manin takara kurus haa maana, ore antee, bean pihi komana.}&When you try to uproot it, it won’t get uprooted, until you fart.\\
\end{xltabular}
\begin{figure}[H]
    \includegraphics[trim=0 0 0 0,clip,width=.5\linewidth]{figures/11_10.jpg}
    \caption{\textit{Kara pihi} -- Farting weed}
    \label{fig:11_10}
\end{figure}
\begin{xltabular}{\textwidth}{lXX}
	(\ref{ex:text7-64})&\textbf{\textit{Kara rovivi}}&\textbf{\textit{Rovivi}, a kind of sappy vine}\\
	(\ref{ex:text7-65})&\textit{O kara rovivi na tei me nana bono marasin.}&The \textit{rovivi} vine contains medicine.\\
	(\ref{ex:text7-66})&\textit{Bean kiu bata nom, ean re mosi vohara a mooman,}&When you work and accidentally cut your leg,\\
	(\ref{ex:text7-67})&\textit{ean toro gono o kara rovivi,}&you must get a \textit{rovivi} vine,\\
	(\ref{ex:text7-68})&\textit{ean re buhaha e.}&then rub it into pulp in your hands.\\
	(\ref{ex:text7-69})&\textit{Ean repaa posie bona mosi.}&Then you squeeze the juice out onto the cut.\\
	(\ref{ex:text7-70})&\textit{Omana a vihii a beera,}&Sure, the stinging pain is great,\\
	(\ref{ex:text7-71})&\textit{evehee pasi too vaha koana a mosi tean.}&but your cut will heal.\\
	(\ref{ex:text7-72})&\textit{Eve a taba to mataa suku nae.}&That's why it is good.\\
	(\ref{ex:text7-73})&\textit{O hum to hata kaakoo nana o kara rovivi, bean mosie, ore paa sapusaka vo tea maa mataman.}&A situation in which the \textit{rovivi} is bad is when you cut it, and the liquid splashes into your eyes.\\
	(\ref{ex:text7-74})&\textit{O kara rovivi mee saka kikis haana.}&Also, the \textit{rovivi} vine is not strong.\\
\end{xltabular}
\begin{figure}[H]
    \includegraphics[trim=0 0 0 3,clip,width=.53\linewidth]{figures/11_11.jpg}\hfil
    \includegraphics[trim=0 0 0 6,clip,width=.47\linewidth]{figures/11_12.jpg}\hfil
    \caption{\textit{Kara rovivi}}
    \label{fig:11_11}
\end{figure}
\begin{figure}[H]
    \captionsetup[subfigure]{oneside,margin={0cm,0cm}}
    \subfloat[]{\includegraphics[trim=0 0 20 2,clip,width=.55\linewidth]{figures/11_13.jpg}\label{fig:11_13}}
    \subfloat[]{\includegraphics[trim=0 0 0 0,clip,width=.45\linewidth]{figures/11_14.jpg}\label{fig:11_14}}
    \caption{(a) \textit{Kara rovivi}, (b) \textit{Buhaha} -- Rub leaves into pulp with your hands}
\end{figure}
\begin{xltabular}{\textwidth}{lXX}
	(\ref{ex:text7-75})&\textbf{\textit{Kara raviana}}&\textbf{\textit{Raviana} vine}\\
	(\ref{ex:text7-76})&\textit{O kara raviana, eve o kara korisi to pampam sinivi ni riori.}&The \textit{raviana} vine, it is a vine that they use to tie up the canoe.\\
	(\ref{ex:text7-77})&\textit{Nabunuu komana vai roho suguna vai nomana.}&In the past until today.\\
	(\ref{ex:text7-78})&\textit{O kara koris bona o kara kikis kurus,}&This vine is very strong,\\
	(\ref{ex:text7-79})&\textit{ei to rake suku riori tea pam sinivi.}&this is why they want (it) for tying up canoes.\footnote{The \textit{raviana} vine is also used for the \textit{kakavei} `bamboo frame net', see Chapter \ref{chap:29}:\ref{ex:text22-63}ff.}\\
\end{xltabular}
\begin{figure}[H]
    \captionsetup[subfigure]{oneside,margin={0cm,0cm}}
    \subfloat[]{\includegraphics[trim=0 0 0 0,clip,width=.47\linewidth]{figures/11_15.jpg}\label{fig:11_15}}
    \subfloat[]{\includegraphics[trim=0 0 0 0,clip,width=.53\linewidth]{figures/11_16.png}\label{fig:11_16}}
    \caption{(a) \textit{Raviana} -- A vine, (b) \textit{Kara raviana} -- \textit{Raviana} string}
\end{figure}
\begin{xltabular}{\textwidth}{lXX}
	(\ref{ex:text7-80})&\textbf{\textit{Kara ruhu}}&\textbf{Aerial roots of the Bengal ficus}\\
	(\ref{ex:text7-81})&\textit{O ruhu eve o meho naono to sun nana rasuu.}&The Bengal ficus is a tree that stands in the jungle.\\
	(\ref{ex:text7-82})&\textit{O naono bona o beera sana.}&This tree is very big.\\
	(\ref{ex:text7-83})&\textit{Eve he o kara berobero.}&It has many aerial roots.\\
	(\ref{ex:text7-84})&\textit{O kara nae na ikuu gunaha taneo maana gina tea maa pea nae.}&Its aerial roots hang down from its branches high up.\\
	(\ref{ex:text7-85})&\textit{O kara ruhu na topo nana gina ore paa hana gunaha maana vo putaa.}&The aerial roots of the Bengal ficus stick there high up and hang down to the ground.\\
	(\ref{ex:text7-86})&\textit{O ana ruhu mee na tei me nana bono porete vaa tea kamis nihi.}&The root of the Bengal ficus contains medicine against tooth ache.\\
	(\ref{ex:text7-87})&\textit{Koto vamuramuraka o sii ana ruhu vai o sii rutaa, ean repaa vaatei e teo hum takepo.}&Bite off and soften a little piece of the Bengal ficus root, then you put it on the place (where a piece of the tooth) split off.\\
\end{xltabular}
\begin{figure}[H]
    \includegraphics[trim=0 0 0 0,clip,width=.85\linewidth]{figures/11_17.jpg}
    \caption{\textit{Ruhu} -- Bengal ficus}
    \label{fig:11_17}
\end{figure}
\begin{xltabular}{\textwidth}{lXX}
	(\ref{ex:text7-88})&\textbf{\textit{Kara urisi}}&\textbf{Strap of a wild banana sheath}\\
	(\ref{ex:text7-89})&\textit{O kara urisi eve to vateen ni raara.}&The wild banana sheath strap, we use for carrying the backpack.\\
	(\ref{ex:text7-90})&\textit{Eara na gono rara o urisi teo veresun.}&We get the sheath from the wild banana plant.\\
	(\ref{ex:text7-91})&\textit{A veresun na tei rori rasuu.}&The wild banana plants are in the jungle.\\
	(\ref{ex:text7-92})&\textit{Eara na kakapa rara o veresun.}&We pull (the sheath) off the wild banana plant.\\
	(\ref{ex:text7-93})&\textit{Are paa varoava e ore mahaka.}&Then we put it in the sun, so that it becomes dry.\\
	(\ref{ex:text7-94})&\textit{Be mahaka, eara re garuhu e, ei beara viru bata rae.}&When (it) is dry, we wash it while we are rolling it up.\\
	(\ref{ex:text7-95})&\textit{Mahaka vai ei, eara repaa gono tea vateen nie.}&Now (it) is dry, and we get it for carrying the backpack.\footnote{The strap for carrying backpacks can also be made from hibiscus fibres, see Chapter \ref{chap:12}:\ref{ex:text8-8}.}\\
\end{xltabular}
\begin{figure}[H]
    \captionsetup[subfigure]{oneside,margin={0cm,0cm}}
    \subfloat[]{\includegraphics[trim=0 70 0 58,clip,width=.26\linewidth]{figures/11_18.jpg}\label{fig:11_18}}\hfil
    \subfloat[]{\includegraphics[trim=0 50 0 26,clip,width=.26\linewidth]{figures/11_19.jpg}\label{fig:11_19}}\hfil
    \subfloat[]{\includegraphics[trim=0 0 200 0,clip,width=.48\linewidth]{figures/11_20.jpg}\label{fig:11_20}}\hfil
    \caption{(a) \textit{Veresun} -- Wild banana, (b) \textit{Kara urisi} -- Wild banana sheath, (c) \textit{Kakapa o veresun} -- Pull (the sheath) off the wild banana plant}
\end{figure}
\begin{figure}[H]
    \captionsetup[subfigure]{oneside,margin={0cm,0cm}}
    \subfloat[]{\includegraphics[trim=0 0 0 0,clip,width=.65\linewidth]{figures/11_21.jpg}\label{fig:11_21}}\hfil
    \subfloat[]{\includegraphics[trim=0 0 0 6,clip,width=.35\linewidth]{figures/11_22.jpg}\label{fig:11_22}}
    \caption{(a) \textit{Varoava} -- Drying in the sun, (b) \textit{Vateen} -- Carrying the backpack}
\end{figure}
\begin{xltabular}{\textwidth}{lXX}
	(\ref{ex:text7-96})&\textbf{\textit{Vunevune}}&\textbf{Beach morning glory}\\
	(\ref{ex:text7-97})&\textit{O vunevune o kara korisi to hirihiri batana teo kasuana.}&The beach morning glory is a vine that crawls along on the ground.\\
	(\ref{ex:text7-98})&\textit{A vuuaba to vaahito rori teo kara vunevune a Nao Boen.}&The clan who owns the beach morning glory is the \textit{Nao Boen}.\\  
\end{xltabular}
\begin{figure}[H]
    \captionsetup[subfigure]{oneside,margin={0cm,0cm}}
    \subfloat[]{\includegraphics[trim=0 0 0 10,clip,width=.56\linewidth]{figures/11_23.jpg}\label{fig:11_23}}\hfil
    \subfloat[]{\includegraphics[trim=0 0 0 0,clip,width=.44\linewidth]{figures/11_24.jpg}\label{fig:11_24}}\hfil
    \caption{\textit{Vunevune} -- Beach morning glory}
\end{figure}
\end{document}